\chapter{Heavy Metal Testing}

\section*{What Is This Test?}
Heavy metal testing measures toxic metals in blood or urine to identify exposure and guide treatment. The most clinically relevant heavy metals include lead, mercury, arsenic, and cadmium, each with distinct exposure sources, toxicities, and diagnostic thresholds \cite{kosnett2024heavy}. Other metals sometimes tested include chromium, cobalt, nickel, thallium, and uranium. Specimen choice depends on the metal: blood is preferred for lead and mercury; 24-hour urine is preferred for arsenic and cadmium. Hair analysis is generally not recommended for clinical diagnosis.

\section*{Alternative Names}
Heavy Metal Panel; Lead Level; Mercury Level; Arsenic Level; Cadmium Level; Toxic Metal Screen; Heavy Metal Blood Test; Heavy Metal Urine Test

\section*{Principle}
Detection is by inductively coupled plasma mass spectrometry (ICP-MS), the gold standard for trace and toxic element analysis \cite{nordberg2022handbook}. Alternative methods include atomic absorption spectroscopy (AAS), especially graphite furnace AAS for lead. Each element has a unique mass-to-charge ratio detected by the mass spectrometer, providing parts-per-billion sensitivity. Whole blood is collected in metal-free EDTA tubes (royal blue-top); urine is collected in acid-washed 24-hour containers.

\section*{History}
Concerns about heavy metal toxicity date to antiquity (lead poisoning in Roman aqueducts, mercury in felters). Occupational exposure limits were established in the 20th century \cite{osha2024heavy}. The U.S. CDC lowered the reference value for lead in children from 10 $\mu$g/dL (2012) to 3.5 $\mu$g/dL (2021), reflecting accumulating evidence of neurocognitive harm at low levels \cite{cdc2021bloodlead}.

\section*{Why This Test Is Ordered}
\begin{itemize}
  \item Acute or chronic symptoms consistent with toxicity (abdominal pain, neurologic dysfunction, encephalopathies, neuropathy, renal failure)
  \item Occupational monitoring of exposed workers (battery, smelting, welding, mining, ammunition, dental amalgam)
  \item Pediatric screening for lead exposure (Medicaid-eligible children at ages 12 and 24 months; high-risk zip codes)
  \item Pre- and post-chelation therapy monitoring
  \item Unexplained chronic kidney disease, peripheral neuropathy, or anemia with basophilic stippling
  \item Suspected suicidal or homicidal ingestion
\end{itemize}

\section*{Primary vs Secondary Test}
Primary for suspected exposure; secondary for unexplained multi-system symptoms.

\section*{How This Test Is Performed}
\textbf{Blood:} Royal blue-top EDTA tube (trace-metal-free) --- venipuncture using stainless steel needle. \textbf{Urine:} 24-hour collection in acid-washed container; a random urine with creatinine correction can be substituted. \textbf{Hair:} not recommended. For arsenic, distinguish toxic inorganic arsenic from benign organic arsenic (seafood) by asking the patient to abstain from seafood for 5 days before testing.

\section*{What Preparation Is Needed}
Avoid seafood for 5 days (for arsenic testing). Avoid chelation therapy for at least 48 hours unless measuring therapeutic response. Document occupational and residential exposures.

\section*{Contraindications and Precautions}
Specimens easily contaminated by collection equipment; only royal blue-top trace-element tubes are acceptable for blood. Glass containers may leach contaminants. Hair analysis has no clinical role for toxicity assessment; alternative practitioners use it but is not regulated.

\section*{Risks}
Venipuncture only.

\section*{What Happens After the Test}
Resultable in 3--7 days. A confirmed elevated result prompts source identification, environmental history, public health notification (for lead in children), and chelation when clinically indicated.

\section*{Understanding Results}

\subsection*{Below Normal}
A normal result reduces concern for chronic toxicity but does not exclude past exposure or body burden (metals accumulate in bone, kidney, brain). Hair analysis unreliable.

\subsection*{Above Normal}
\begin{itemize}
  \item \textbf{Lead:} Encephalopathy, peripheral neuropathy, abdominal colic, microcytic anemia with basophilic stippling, nephropathy, hypertension. Children: lowered IQ, behavioral issues at levels $>$ 3.5 $\mu$g/dL
  \item \textbf{Mercury:} Neurologic (tremor, memory, ataxia), renal tubular damage, gingivitis, acrodynia in children. Forms: elemental (vapor), inorganic ( salts ), organic (methylmercury in fish)
  \item \textbf{Arsenic:} Acute: GI distress, cardiac arrhythmia, encephalopathy, hemolysis. Chronic: skin (hyperkeratosis, hyperpigmentation), neuropathy, peripheral vascular disease, cancers (skin, lung, bladder)
  \item \textbf{Cadmium:} Renal tubular dysfunction (proteinuria), osteomalacia, pulmonary disease, increased cancer risk
  \item \textbf{Other metals:} Cobalt (cardiomyopathy), thallium (alopecia, neuropathy), chromium (renal, dermal)
\end{itemize}

\section*{Normal Results}
\begin{itemize}
  \item \textbf{Lead (whole blood):} Adult $<$ 5 $\mu$g/dL; Child $<$ 3.5 $\mu$g/dL (CDC reference value); action levels: 3.5--45 $\mu$g/dL monitor and reduce sources; $>$ 45 chelation; $>$ 70 urgent chelation
  \item \textbf{Mercury (whole blood):} $<$ 5 $\mu$g/L; toxicity $>$ 15 $\mu$g/L
  \item \textbf{Inorganic arsenic (24-h urine):} $<$ 20 $\mu$g/L (seafood refrain needed)
  \item \textbf{Cadmium (blood):} $<$ 5 $\mu$g/L; (urine): $<$ 2 $\mu$g/g creatinine
  \item \textbf{Urine mercury:} $<$ 10 $\mu$g/L
\end{itemize}

\section*{Follow-up Tests}
\begin{itemize}
  \item \textbf{CBC with peripheral smear:} Basophilic stippling in lead poisoning; hemolytic anemia in arsenic
  \item \textbf{Renal and liver function tests:} Target organ effects
  \item \textbf{Iron studies and ZPP (zinc protoporphyrin):} Chronic lead exposure
  \item \textbf{Nerve conduction studies:} Symptomatic peripheral neuropathy
  \item \textbf{Environmental investigation:} Public health testing of home, water, workplace, toys
  \item \textbf{Chelation therapy and post-chelation monitoring:} Succimer (DMSA), EDTA, BAL (dimercaprol), D-penicillamine, depending on metal \cite{acmt2010position}
  \item \textbf{Pregnancy and breastfeeding status:} Lead readily crosses placenta
\end{itemize}

\section*{References}
\bibliographystyle{unsrtnat}
\bibliography{references}

\clearpage
\section*{Regulatory Status and Authorization Date}
Regulatory status applies to the specific assay, instrument, manufacturer, and intended use rather than to the test name alone. No single approval date can be assigned to this test category. For a commercial assay, verify the current product in the relevant regulator's database: the U.S. Food and Drug Administration (FDA) clearance, approval, or authorization; India's Central Drugs Standard Control Organisation (CDSCO) licence or permission; the European Union In Vitro Diagnostic Regulation (EU IVDR) CE certificate issued through the applicable conformity-assessment route; or the United Kingdom Medicines and Healthcare products Regulatory Agency (MHRA) registration and UKCA/CE status. Laboratory-developed tests may not have an individual FDA, CDSCO, EU IVDR, or MHRA approval date.
\clearpage
